\subsection{APPLICATIONS OF CONJUGACY}

THEOREM 2. Let \(\mathfrak{g}\) be a Lie algebra.

\begin{enumerate}
\def\labelenumi{(\roman{enumi})}
\item
  The Cartan subalgebras of \(\mathfrak{g}\) are all of the same dimension, namely \(\operatorname{rk}(\mathfrak{g})\) , and the same nilpotency class.
\item
  An element \(x \in \mathfrak{g}\) is regular if and only if \(\mathfrak{g}^0(x)\) is a Cartan subalgebra of \(\mathfrak{g}\) ; every Cartan subalgebra is obtained in this way.
\end{enumerate}

To prove (i), we can assume that k is algebraically closed (cf.~§2, Prop. 3 and Prop. 6), in which case it follows from Th. 1 of no. 2. Assertion (ii) follows from (i) and §2, Th. 1 (i) and (iv).

PROPOSITION 3. Let \(\mathfrak{g}\) be a Lie algebra, \(\mathfrak{g}'\) a subalgebra of \(\mathfrak{g}\) . The following conditions are equivalent:

\begin{enumerate}
\def\labelenumi{(\roman{enumi})}
\item
  \(\mathfrak{g}'\) contains a regular element of \(\mathfrak{g}\) , and \(\mathrm{rk}(\mathfrak{g}) = \mathrm{rk}(\mathfrak{g}')\) ;
\item
  \(\mathfrak{g}'\) contains a Cartan subalgebra of \(\mathfrak{g}\) ;
\item
  every Cartan subalgebra of \(\mathfrak{g}'\) is a Cartan subalgebra of \(\mathfrak{g}\) .
\item
  \(\Longrightarrow\) (ii): Assume that \(\operatorname{rk}(\mathfrak{g}) = \operatorname{rk}(\mathfrak{g}')\) , and that there exists \(x \in \mathfrak{g}'\) regular in \(\mathfrak{g}\) . Put \(\mathfrak{h} = \mathfrak{g}^0(x)\) , \(\mathfrak{h}' = \mathfrak{g}'^0(x) = \mathfrak{h} \cap \mathfrak{g}'\) . Then
\end{enumerate}

\[
\operatorname{rk} \left(\mathfrak {g} ^ {\prime}\right) \leq \dim \mathfrak {h} ^ {\prime} \leq \dim \mathfrak {h} = \operatorname{rk} (\mathfrak {g}) = \operatorname{rk} \left(\mathfrak {g} ^ {\prime}\right)
\]

so \(\mathfrak{h} = \mathfrak{h}'\subset \mathfrak{g}'\) . This proves (ii).

\begin{enumerate}
\def\labelenumi{(\roman{enumi})}
\setcounter{enumi}{1}
\item
  \(\Longrightarrow\) (iii): Assume that \(g'\) contains a Cartan subalgebra h of g, and let \(h_{1}\) be a Cartan subalgebra of \(g'\) . To prove that \(h_{1}\) is a Cartan subalgebra of g, we can assume that k is algebraically closed. Let \(\mathrm{E}(\mathfrak{h})\) and \(\mathrm{E}'(\mathfrak{h})\) be the groups of automorphisms of g and \(g'\) associated to h (no. 2). By Th. 1, there exists \(f \in \mathrm{E}'(\mathfrak{h})\) such that \(f(\mathfrak{h}) = \mathfrak{h}_{1}\) . Now every element of \(\mathrm{E}'(\mathfrak{h})\) is induced by an element of \(\mathrm{E}(\mathfrak{h})\) ; indeed, it suffices to verify this for \(e^{\operatorname{ad} x}\) , with \(x \in \mathfrak{g}'^{\lambda}(\mathfrak{h}), \lambda \neq 0\) , in which case it follows from the inclusion \(\mathfrak{g}'^{\lambda}(\mathfrak{h}) \subset \mathfrak{g}^{\lambda}(\mathfrak{h})\) . Thus \(h_{1}\) is a Cartan subalgebra of g.
\item
  \(\Longrightarrow\) (i): Assume that condition (iii) is satisfied. Let h be a Cartan subalgebra of \(g'\) . Since this is a Cartan subalgebra of g, it contains a regular element of g (Th. 2 (ii)), and on the other hand \(\operatorname{rk}(\mathfrak{g}) = \dim(\mathfrak{h}) = \operatorname{rk}(\mathfrak{g}')\) .
\end{enumerate}

COROLLARY. Let \(\mathfrak{h}\) be a nilpotent subalgebra of \(\mathfrak{g}\) . The subalgebra \(\mathfrak{g}^0 (\mathfrak{h})\) has properties (i), (ii), (iii) in Prop. 3.

Indeed, Prop. 11 of §2, no. 3, shows that \(\mathfrak{g}^0 (\mathfrak{h})\) has property (ii).

PROPOSITION 4. Let g be a Lie algebra, l the rank of g, c the nilpotency class of the Cartan subalgebras of g, and \(x \in g\) . There exists an l-dimensional subalgebra of g whose nilpotency class is \(\leq c\) and which contains x.

Let T be an indeterminate. Let \(k' = k(T)\) and \(g' = g \otimes_k k'\) . If h is a Cartan subalgebra of g, \(h \otimes_k k'\) is a Cartan subalgebra of \(g'\) , hence the rank of \(g'\) is l and the nilpotency class of the Cartan subalgebras of \(g'\) is c.

Choose a regular element \(y\) of \(\mathfrak{g}\) . With the notations of §2, no. 2, we have \(a_l(y) \neq 0\) . Denote also by \(a_l\) the polynomial function on \(\mathfrak{g}'\) that extends \(a_l\) . Then the element \(a_l(x + \mathrm{T}y)\) of \(k[\mathrm{T}]\) has dominant coefficient \(a_l(y)\) . In particular, \(x + \mathrm{T}y\) is regular in \(\mathfrak{g}'\) . Let \(\mathfrak{h}'\) be the nilspace of \(\operatorname{ad}(x + \mathrm{T}y)\) in \(\mathfrak{g}'\) . Then \(\dim \mathfrak{h}' = l\) and the nilpotency class of \(\mathfrak{h}'\) is \(c\) . Put \(\mathfrak{k} = \mathfrak{h}' \cap (\mathfrak{g} \otimes_k k[\mathrm{T}])\) ; then \(\mathfrak{k} \otimes_{k[\mathrm{T}]} k(\mathrm{T}) = \mathfrak{h}'\) .

Let \(\varphi\) be the homomorphism from \(k[\mathrm{T}]\) to \(k\) such that \(\varphi (\mathrm{T}) = 0\) , and let \(\psi\) be the homomorphism \(1\otimes \varphi\) from \(\mathfrak{g}\otimes_k k[\mathrm{T}]\) to \(\mathfrak{g}\) . Then \(\psi (\mathfrak{k})\) is a subalgebra of \(\mathfrak{g}\) whose nilpotency class is \(\leq c\) and which contains \(\psi (x + \mathrm{T}y) = x\) .

In the free \(k[\mathrm{T}]\) -module \(\mathfrak{g} \otimes_k k[\mathrm{T}]\) , \(\mathfrak{k}\) is a submodule of rank \(l\) , and \((\mathfrak{g} \otimes_k k[\mathrm{T}]) / \mathfrak{k}\) is torsion free, so the submodule \(\mathfrak{k}\) is a direct summand of \(\mathfrak{g} \otimes_k k[\mathrm{T}]\) (Algebra, Chap. VII, §4, no. 2, Th. 1). Hence \(\dim_k \psi(\mathfrak{k}) = l\) , which completes the proof.
% semantic-fragments: legacy-input-seam
\relax{}
